\begin{frame}[plain]
\begin{center}
  \includegraphics[width=0.88\linewidth]{sigmoid_relu_gradient.png}
  \vskip0.4em
  {\footnotesize 시그모이드는 $|z|$가 크면 도함수가 0에 가까워져
  사라지는 그래디언트가 생기지만, ReLU는 $z>0$에서 도함수가 항상
  1 --- 이번 학기 실습의 기본 활성 함수로 ReLU를 쓰는 이유.}
\end{center}
\end{frame}
